\documentclass[10pt]{article}
\usepackage[margin=20mm]{geometry}
\usepackage{amsmath,amssymb,mathtools,fontspec,unicode-math,xurl,hyperref}
\IfFontExistsTF{DejaVu Serif}{\setmainfont{DejaVu Serif}}{\setmainfont{Latin Modern Roman}}
\setmathfont{latinmodern-math.otf}
\setlength{\parindent}{0pt}\setlength{\parskip}{6pt}\sloppy
\title{Probabilistic Construction of Kakeya-Type Sets in the Plane Associated to Separated Directions:
Review of the Acquired Sources}
\author{AI-assisted technical review}
\date{}
\begin{document}
\maketitle
Abstract: The fixed-separation theorem is supported by a coherent pruning, geometric and percolation argument. A separate construction in the concluding section does not, however, define the claimed infinite-lacunary direction set: with the printed order of the interval operators, the set is empty. This is a Major defect in that substantive example, with a precise correction and a proof of the intended separation limitation. One Minor correction concerns an impossible area comparison. It does not refute the main theorem; a uniform rectangular replacement is supplied. The displayed sticky map satisfies the required prefix conditions in both acquired PDFs and the obtained TeX. It also makes the independent choices in the percolation argument explicit. The author-hosted journal-layout reproduction is distinguished from publisher-original bytes and from the separately obtained preprint. Adjudication remains pending.\par
\section*{Sources and the scope of the conclusions}
This is the technical review of M25-0198: Paul Hagelstein, Blanca Radillo-Murguia and Alexander Stokolos, Probabilistic construction of Kakeya-type sets in \(\mathbb R^2\) associated to separated sets of directions [primary], DOI https://doi.org/10.1215/00127094-2025-0025. The primary is a 14-page author-hosted reproduction with journal pagination 185–198, Duke Mathematical Journal 175(2). It was retrieved from the author’s Baylor site; its Acrobat Distiller metadata dates the reproduction to 26 February 2026. Publisher-original byte identity has not been established. Its text layer is garbled, so the primary review used all fourteen actual page images.\par
The first page states first online publication on 21 December 2025. The frozen batch’s OpenAlex membership date, 22 December 2025, and the acquired Crossref print date, 1 February 2026, are distinct metadata fields. The separate arXiv:2405.17674v2 PDF, dated 10 July 2025, has thirteen pages [preprint]. The obtained TeX was fully read without execution. The two Table 1 entries checked directly in both acquired PDFs agree with that TeX. This local agreement does not establish global equivalence between the versions. The 25 manifest-bound source files total 4,069,508 bytes, with no intake size or hash anomaly. No separate supplement or released code was located in the bounded acquisition searches; the publisher supplement inventory was inaccessible. This is a source limitation, not a mathematical defect.\par
Theorem 1, primary printed page 189 (PDF page 5), assumes \[\exists\,\eta>0\quad\forall\,N\in\mathbb N:
 \quad \mathcal T_\Omega\text{ contains an }\eta\text{-separated
 subtree of lacunary order }N.
 \tag{1}\] The same \(\eta\) must work for all \(N\). The conclusion is unboundedness of the rectangular directional maximal operator on every \(L^p(\mathbb R^2)\), \(1\le p<\infty\). It does not assert unboundedness on \(L^\infty\), nor does it replace (1) by infinite lacunary value alone.\par
The tree uses closed dyadic intervals; a dyadic endpoint may give two rays. The manuscript explicitly acknowledges this and distinguishes its lacunary value from conventional lacunary order. Separation concerns two splitting descendants on opposite halves of a splitting ancestor, with distance at least \(\eta\) times the ancestor’s length. It is not a uniform separation of every pair of leaves.\par
The proof prunes to \(2^N\) terminal vertices at a common finite height and forms strips from sticky maps. Lemma 1 provides a near-slab area lower bound \(\gtrsim(\log N)/N\). Lemma 2 supplies a map with far-slab area \(\lesssim_\eta 1/N\). A directional maximal estimate then turns their ratio into unbounded operator norms. Theorem 2 separately gives a probability-one counterexample to a pointwise step in an earlier argument. The concluding section promises an infinite-lacunary set outside (1); that promise is the locus of the Major issue.\par
\section*{The concluding direction set is empty as printed}
Problem 1 (Section 5, printed page 197; PDF page 13).  For an interval \(I=[a,b]\), the source defines \[\rho_{j,1}(I)=
 \left[\frac{a+b}{2}-2^{-j}(b-a),\frac{a+b}{2}\right],
 \qquad
 \rho_{j,2}(I)=
 \left[\frac{a+b}{2},\frac{a+b}{2}+2^{-j}(b-a)\right].
 \tag{2}\] It defines \(\Omega\) by membership in infinitely many intervals \[\rho_{1,i_1}\rho_{2,i_2}\cdots\rho_{k,i_k}[0,1],
 \qquad i_j\in\{1,2\},
 \tag{3}\] and claims \(\lambda(\Omega)=\infty\). Under the composition convention used in the preceding section, (3) instead gives \(\Omega=\varnothing\).\par
These are operators on intervals, depending on each current midpoint and length. They cannot be replaced by fixed affine maps on points while retaining (3). Each \(\rho_{j,i}\) sends an interval into itself. Consequently every interval in the union \(U_k\) at level \(k\) is contained in the first interval acted upon: \[U_k\subset
 \left[\frac12-2^{-k},\frac12+2^{-k}\right].
 \tag{4}\] For \(k\ge3\), put \(\varepsilon=2^{-k}\). The innermost image is either \([1/2-\varepsilon,1/2]\) or \([1/2,1/2+\varepsilon]\). The next operation \(\rho_{k-1,i}\) places its image strictly inside that interval, at distance at least \[\varepsilon\left(\frac12-2^{-(k-1)}\right)>0
 \tag{5}\] from its endpoint \(1/2\). Every subsequent operation stays inside the current interval. Thus \(1/2\notin U_k\) for all \(k\ge3\).\par
Every fixed \(x\ne1/2\) is excluded from \(U_k\) for all large \(k\) by (4); the only possible remaining point is excluded by (5). There are only finitely many intervals at any fixed level, so no point belongs to infinitely many intervals in (3). This proves the assertion of emptiness, without a numerical or asymptotic assumption. An exact finite check gives \[U_3=[13/32,15/32]\cup[17/32,19/32].\] The source initially defines its directional objects for nonempty sets; on that convention the proposed example is inadmissible. Under the natural extension to the root-only tree, its splitting value is zero, never infinity.\par
\subsection*{Counterevidence and impact}
A left-to-right application convention would produce the intended forward recursion, but the source does not specify one. More decisively, Section 4, printed page 195, explicitly writes \(\rho_{a_k}\cdots\rho_{a_1}I_0\) for its forward recursive construction. The same ordinary convention makes (3) the reversed construction. Both obtained PDFs and the obtained TeX retain this discrepancy.\par
This is Major for the claimed Section 5 example: the construction supplies no infinite-lacunary direction set, and its stated model question concerns the wrong set. The example is a substantive claimed distinction, announced in the introduction, between infinite splitting value and the stronger hypothesis (1). Its failure is not merely an omitted explanatory sentence. The defect does not refute Theorem 1, Theorem 2, or the possibility of the intended example. Its severity is based on the failed mathematical conclusion, despite the short syntactic correction.\par
\subsection*{A proved repair, including arbitrary subtrees}
Replace (3) by \[I_{i_1\ldots i_k}
 =\rho_{k,i_k}\cdots\rho_{1,i_1}[0,1].
 \tag{6}\] The nested construction intervals have lengths \(L_k=2^{-k(k+1)/2}\). In binary notation the left child at stage \(k\) appends a block \(0\,1^{k-1}\), and the right child appends \(1\,0^{k-1}\). Every block for \(k\ge2\) contains both digits. Hence no infinite address is eventually constant and no limit point is dyadic. The closed-endpoint convention introduces no extra rays into this corrected example.\par
The nested compact unions define a nonempty compact set \(C\). Its splitting vertices are precisely the construction intervals, with unary paths between construction generations. Every finite number of generations gives a finite binary branching pattern, so \(\lambda(C)=\infty\).\par
Fix \(\eta>0\), and choose \(J\) with \(2^{-J}<\eta\). For a construction interval \(I\) of generation \(j\ge J\), all its points of \(C\) lie in the central union of its two children, whose total span is \[2L_{j+1}=2^{-j}L_j<\eta |I|.\] If a subtree has splitting descendants \(u,v\) in opposite halves of \(I\), their construction intervals satisfy \(\operatorname{dist}(Q_u,Q_v)<\eta |I|\), contrary to separation. Therefore, in an \(\eta\)-separated subtree, all splitting vertices in each such late descendant region lie along one ray. Otherwise their least common ancestor would violate the preceding condition. Each tail has lacunary order at most one. The first \(J\) construction generations add at most \(J\) to that order, giving an upper bound \(J+1\) for the whole subtree. This includes subtrees that skip generations. Thus (6) has exactly the intended property, and the repair does not rely on an unproved assertion about a selected pruning.\par
\section*{A local correction with preserved main estimates}
Problem 2 (Reciprocal area, printed page 192; PDF page 8).  Put \(w=2^{-h(\mathcal P)}\) and \(L=(2\eta)^{-1}+1\). The proof says that a parallelogram is contained in a direction-adjusted parallelogram with \(1/L\) times its area. Since \(L>1\), positive-area containment makes that assertion impossible.\par
For example, \(\eta=1/4\) gives \(L=3\). The original parallelogram has area \(3w\); the alleged containing one would have area \(w\). This direct contradiction applies to the actual geometric claim. The preceding average identity is nevertheless correct: the far set is restricted to a unit-width \(x\)-slab, so its intersection with the given strip has area \(w\), and its average over the length-\(L\) parallelogram is \(1/L\).\par
Here is a full replacement for the required maximal estimate. Choose \(a\in\Omega\) in the terminal dyadic interval whose left endpoint is \(d\); thus \(0\le a-d\le w\). For the parallelogram \[P=\{(x,y):0\le x\le L,\ b+dx<y<b+dx+w\},\] its orthogonal span transverse to \((1,a)\) is at most \((L+1)w\), and its longitudinal span is at most \(\sqrt2L+w\le3L\). For the theorem we may decrease the fixed separation parameter to ensure \(\eta\le1/2\), hence \(L\ge2\). Taking a containing rectangle with long edge parallel to \((1,a)\), and padding to make it open, gives \[|R|\le6L(L+1)w,\qquad
 \int_R\chi_{K_{\sigma,2,\eta}}\ge w.\] The padding can be chosen within the stated factor. Every point of the near strip is then in an admissible rectangle with average at least \[\frac1{6L(L+1)}\gtrsim\eta^2,
 \tag{7}\] uniformly in \(N\) and the terminal height. Together with the two area estimates, for each finite \(p\ge1\) this gives \[\|M_\Omega\|_{L^p\to L^p}
 \ge c\eta^2
 \left(\frac{|K_{\sigma,1}|}{|K_{\sigma,2,\eta}|}\right)^{1/p}
 \gtrsim_{\eta,p}(\log N)^{1/p}\longrightarrow\infty.\] Thus the main operator estimate survives with a correct enlargement. This is Minor: the numerical area assertion is false, but the needed conclusion follows by the explicit uniform replacement above. It is not classified as Minor merely because one symbol can be changed.\par
\subsection*{The displayed sticky map}
Table 1, printed page 190 (PDF page 6), lists \(0101\mapsto1010\) and \(0110\mapsto1010\). Both entries are also present in the separate arXiv PDF and the obtained TeX. All sixteen displayed images preserve the common prefixes of their inputs, so the table does define a sticky map. The earlier alleged counterexample used incorrectly transcribed output strings and is withdrawn; no change to the authors’ table is indicated.\par
\section*{Positive checks and proof dependencies}
The separation estimate in Section 3 has the correct geometric scale. For a splitting interval \(I\) and two opposite splitting descendants, admissible slopes differ by at least \(\eta |I|\). At \(x>(2\eta)^{-1}\) their intercepts differ by more than \(|I|/2\), so they cannot both occur within the same half of a base interval of length \(|I|\). The restriction on the vertical boundary is harmless for the area integral.\par
For precision, the random sticky map can be generated using independent fair choices at input edges whenever the current target vertex splits. The phrase over all sticky maps should not be used to assume that uniform counting of the finite set of completed maps has that law. Completion counts can differ when target branching heights differ. Existence in Lemma 2 only requires the fair generative law, which is available.\par
Several candidate intervals inside the same input half share a single choice. They should be grouped under that half, rather than assigned fictitiously independent bits. The compatible half-intervals, joined by containment and with a virtual root, give a genuine independent percolation tree. If \[M_\eta=\left\lceil (2\eta)^{-1}+1\right\rceil+2,\] its level \(d\) contains at most \(M_\eta2^d\) vertices. This direct count avoids changing the far slab merely to make \(\eta\) dyadic.\par
For completeness, the survival bound can be checked elementarily. Give edges entering level \(d\) resistance \(2^{d-1}\). Every unit flow to level \(n\), by Cauchy–Schwarz on each level, has energy at least \(n/(2M_\eta)\). Thus the effective conductance \(C\le2M_\eta/n\). For independent retention probability \(1/2\), induction with \(\varrho(q)=q/(2-q)\) gives \(\varrho(q)\le C\): at a vertex conductances combine as \(\sum_i C_i/(C_i+2)\), and \[\varrho(a+b-ab)
 =\frac{\varrho(a)+\varrho(b)}
        {1+\varrho(a)\varrho(b)}
 \le\varrho(a)+\varrho(b).\] Consequently \(q\le4M_\eta/n\). Taking \(n=N-1\) and handling bounded \(N\) directly gives the required \(O_\eta(N^{-1})\) estimate. On the far slab the union lies in \(0\le y\le1+x\), a region of finite area depending only on \(\eta\). Integration and expectation therefore provide the map claimed in Lemma 2. These completions preserve the actual target estimate; no counterexample to it was established.\par
Theorem 2 was checked at its exact level of assertion. Its intervals have length \(8^{-N}\), the offset \(2^{-100N}\) selects the intended finite dyadic tree, and the translated reflected child intervals lie in the same input half. The recursive choice therefore produces a strip containing the displayed point for every sticky map. The point lies \(8^{-(N+1)}\) below its top edge, strictly inside it. This disproves a universal pointwise probability bound. A single point of probability one does not, by itself, disprove the integral Claim 7(B) or the earlier final classification theorem. The source itself makes this distinction; the review retains it. The strict margin also gives a small common region at \(x>1\): if \(\varepsilon=8^{-(N+1)}\), then \(0\le x-1\le\varepsilon/4\) and \(|y-y_N|\le\varepsilon/4\) stay in every \(K_\sigma\). Its area is only \(\varepsilon^2/8\), so this strengthening still gives no contradiction to the integral bound.\par
The constant-ratio finite examples themselves contain separated subtrees after skipping generations: within an interval of length \(L\), the descendants labelled \(11\) and \(22\) are \([27/64,28/64]\) and \([36/64,37/64]\) in relative coordinates, with gap \(L/8\). Recursing on them verifies the opening remark of Section 5. This differs from the variable-ratio construction (2). The final question also differs between acquired versions: the primary asks about boundedness for some \(1<p<\infty\), whereas the preprint asks about every \(1<p\le\infty\). The review does not silently equate those formulations.\par
\section*{Assessment and verification limits}
The confirmed issues are one Major failure of the Section 5 example and one Minor defect in an intermediate area assertion. No defect is established in the illustrative table. The corrected direction set has the claimed separation limitation, and (7) restores the main directional estimate. The main fixed-separation theorem is not disproved by these findings. Nor does the review settle the concluding boundedness question for the corrected set.\par
The inherited near-slab lower bound and pruning are distinct dependencies from the disputed earlier pointwise upper bound. Their actual source and scope are recorded in the supporting dependency note; a separately retrieved 2007 preprint [bateman] is not represented as byte-identical to the 2009 journal article. Claims concerning all downstream applications or the current status of the broader classification theorem are outside this bounded review.\par
Three native helpers supplied bounded tree, probability and construction checks, with separate source and image hash receipts. They are assistance in the same initial review, not independent human acceptance. No author code or author TeX was executed. The delivery record retains compiler dependencies, failed or preliminary records and original render-only flags. The scientific decision remains pending, with an empty human-review list.\par
9 P. Hagelstein, B. Radillo-Murguia, and A. Stokolos, Probabilistic construction of Kakeya-type sets in \(\mathbb R^2\) associated to separated sets of directions, author-hosted journal-layout reproduction, Duke Math. J. 175(2), 185–198, https://bpb-us-w2.wpmucdn.com/sites.baylor.edu/dist/f/83/files/2026/02/Duke-2026.pdf. Primary PDF SHA-256: `638dac82f914f7fdc3ffa304b803fa1d99e499299da2855572ec4f3bfa35b18d`.\par
The same authors and title, arXiv:2405.17674v2, 10 July 2025, https://arxiv.org/abs/2405.17674v2. Separately obtained PDF and TeX; no global equivalence claim is made.\par
M. Bateman, Kakeya sets and directional maximal operators in the plane, arXiv:math/0703559v1, 19 March 2007, https://arxiv.org/abs/math/0703559v1. Bounded dependency check only; distinct from the cited Duke Math. J. 147 (2009), 55–77 journal bytes.\par
\end{document}
